\documentclass[11pt]{article}

% ------------------------------------------------------------
% PACKAGES
% ------------------------------------------------------------
\usepackage[margin=1in]{geometry}
\usepackage{amsmath,amssymb,amsfonts}
\usepackage{booktabs}
\usepackage{array}
\usepackage{longtable}
\usepackage{enumitem}
\usepackage{xcolor}
\usepackage{hyperref}
\usepackage{titlesec}
\usepackage{fancyhdr}
\usepackage{multicol}

% ------------------------------------------------------------
% DOCUMENT SETTINGS
% ------------------------------------------------------------
\hypersetup{
    colorlinks=true,
    linkcolor=blue,
    urlcolor=blue,
    citecolor=blue
}

\setlist[itemize]{leftmargin=1.5em}
\setlist[enumerate]{leftmargin=1.8em}

\pagestyle{fancy}
\fancyhf{}
\rhead{Calculus 1 for Biological and Social Sciences}
\lhead{Full Course Breakdown}
\cfoot{\thepage}

\titleformat{\section}
  {\Large\bfseries}
  {\thesection}
  {0.75em}
  {}

\titleformat{\subsection}
  {\large\bfseries}
  {\thesubsection}
  {0.75em}
  {}

% ------------------------------------------------------------
% CUSTOM COMMANDS
% ------------------------------------------------------------
\newcommand{\objective}[1]{\item \textbf{Learning Objective:} #1}
\newcommand{\concept}[1]{\item \textbf{Concepts:} #1}
\newcommand{\difficult}[1]{\item \textbf{Difficult Topics:} #1}
\newcommand{\prereq}[1]{\item \textbf{Prerequisites:} #1}

% ------------------------------------------------------------
% TITLE
% ------------------------------------------------------------
\begin{document}

\begin{titlepage}
    \centering
    \vspace*{1.5cm}

    {\Huge\bfseries Calculus 1 for Biological and Social Sciences\par}
    \vspace{0.5cm}
    {\LARGE Full Course Structure and Learning Framework\par}

    \vspace{1.5cm}

    {\large
    Chapters, Modules, Concepts, Learning Objectives,\\
    Prerequisites, and Difficult Topics
    \par}

    \vfill

    \textbf{Course Design Basis}\\[0.3cm]
    The structure follows the supplied textbook Table of Contents.
    
    \vspace{0.5cm}
    \textit{Note: Sections marked with an asterisk in the textbook
    are treated as optional or enrichment material.}

    \vfill

    {\large \today\par}
\end{titlepage}

\tableofcontents
\newpage

% ============================================================
% 1. COURSE OVERVIEW
% ============================================================

\section{Course Overview}

This course develops calculus and related mathematical tools for
applications in the biological and social sciences. The progression
moves from mathematical preparation and elementary functions to
limits, differentiation, integration, differential equations,
linear algebra, multivariable calculus, and probability and statistics.

The central conceptual progression is:

\[
\boxed{
\text{Functions}
\rightarrow
\text{Limits}
\rightarrow
\text{Derivatives}
\rightarrow
\text{Applications}
\rightarrow
\text{Integrals}
\rightarrow
\text{Differential Equations}
}
\]

The later portion of the textbook extends this foundation to

\[
\boxed{
\text{Linear Algebra}
\rightarrow
\text{Multivariable Calculus}
\rightarrow
\text{Systems of Differential Equations}
\rightarrow
\text{Probability and Statistics}.
}
\]

\subsection{Primary Course Goals}

By the end of the full course, students should be able to:

\begin{enumerate}
    \item Work fluently with algebraic, graphical, exponential,
    logarithmic, and trigonometric functions.
    
    \item Construct and interpret mathematical models of biological
    and social phenomena.
    
    \item Understand the conceptual meaning of a limit and continuity.
    
    \item Compute and interpret derivatives as rates of change.
    
    \item Use derivatives to analyze functions and solve optimization
    problems.
    
    \item Understand integration as accumulation and inverse
    differentiation.
    
    \item Apply definite and indefinite integrals to scientific models.
    
    \item Solve basic differential equations arising from biological
    and social models.
    
    \item Work with systems of linear equations and matrices.
    
    \item Extend differentiation and modeling ideas to several
    independent variables.
    
    \item Analyze systems of differential equations qualitatively
    and quantitatively.
    
    \item Apply probability and statistical concepts to data and
    stochastic models.
\end{enumerate}

% ============================================================
% 2. PREREQUISITE STRUCTURE
% ============================================================

\section{Prerequisite Structure}

\subsection{Mathematical Prerequisites}

The following skills form the foundation for the course:

\begin{itemize}
    \item Arithmetic and numerical manipulation
    \item Algebraic simplification
    \item Factoring
    \item Solving equations and inequalities
    \item Coordinate geometry
    \item Slopes and equations of lines
    \item Quadratic equations
    \item Basic trigonometry
    \item Exponential functions
    \item Logarithmic functions
    \item Function notation
    \item Graph interpretation
\end{itemize}

\subsection{Prerequisite Dependency Map}

\[
\begin{array}{c}
\text{Algebra + Trigonometry + Functions}
\\[0.3cm]
\downarrow
\\[0.3cm]
\text{Limits and Continuity}
\\[0.3cm]
\downarrow
\\[0.3cm]
\text{Differentiation}
\\[0.3cm]
\downarrow
\\[0.3cm]
\text{Applications of Differentiation}
\\[0.3cm]
\downarrow
\\[0.3cm]
\text{Integration}
\\[0.3cm]
\downarrow
\\[0.3cm]
\text{Integration Techniques}
\\[0.3cm]
\downarrow
\\[0.3cm]
\text{Differential Equations}
\end{array}
\]

The later subjects depend on this core:

\[
\begin{array}{ccccc}
\text{Differential Equations}
&\longrightarrow&
\text{Systems of Differential Equations}
\\
\updownarrow
&&
\updownarrow
\\
\text{Linear Algebra}
&\longrightarrow&
\text{Multivariable Calculus}
\end{array}
\]

Probability and statistics constitute a related mathematical
application stream that can be developed after the foundational
algebra and calculus material.

% ============================================================
% CHAPTER 1
% ============================================================

\section{Chapter 1: Preview and Review}

\subsection{Major Topic}

\textbf{Precalculus skills, elementary functions, and graphing.}

\subsection{Module 1.1: Precalculus Diagnostic}

\textbf{Purpose:} Identify mathematical preparation needed before
formal calculus.

\begin{itemize}
    \concept{Algebraic manipulation}
    \concept{Functions}
    \concept{Equations}
    \concept{Graphs}
    \concept{Trigonometry}
    \concept{Exponentials and logarithms}
\end{itemize}

\subsection{Module 1.2: Preliminaries}

\begin{itemize}
    \concept{Real numbers}
    \concept{Lines in the plane}
    \concept{Equation of a circle}
    \concept{Trigonometry}
    \concept{Exponentials and logarithms}
    \concept{Complex numbers}
    \concept{Quadratic equations}
\end{itemize}

\subsection{Module 1.3: Elementary Functions}

\begin{itemize}
    \concept{Definition of a function}
    \concept{Polynomial functions}
    \concept{Rational functions}
    \concept{Power functions}
    \concept{Exponential functions}
    \concept{Inverse functions}
    \concept{Logarithmic functions}
    \concept{Trigonometric functions}
\end{itemize}

\subsection{Module 1.4: Graphing}

\begin{itemize}
    \concept{Graphing functions}
    \concept{Basic transformations}
    \concept{Logarithmic scales}
    \concept{Transformations into linear functions}
    \concept{Translating verbal descriptions into graphs}
\end{itemize}

\subsection{Learning Objectives}

Students should be able to:

\begin{enumerate}
    \objective{Use real-number and algebraic operations correctly.}
    \objective{Interpret slope and equations of lines.}
    \objective{Use basic trigonometric relationships.}
    \objective{Manipulate exponential and logarithmic expressions.}
    \objective{Define and evaluate functions.}
    \objective{Identify domains and ranges of elementary functions.}
    \objective{Graph polynomial, rational, power, exponential,
    logarithmic, and trigonometric functions.}
    \objective{Apply transformations to function graphs.}
    \objective{Translate between verbal, algebraic, and graphical
    representations.}
\end{enumerate}

\subsection{Prerequisites}

\begin{itemize}
    \prereq{High-school algebra}
    \prereq{Coordinate geometry}
    \prereq{Basic trigonometry}
\end{itemize}

\subsection{Difficult Topics}

\begin{itemize}
    \difficult{Function notation and domain restrictions}
    \difficult{Inverse functions}
    \difficult{Logarithmic transformations}
    \difficult{Graph transformations}
    \difficult{Connecting algebraic equations to graphs}
\end{itemize}

% ============================================================
% CHAPTER 2
% ============================================================

\section{Chapter 2: Discrete-Time Models, Sequences, and Difference Equations}

\subsection{Major Topic}

\textbf{Discrete mathematical modeling and sequences.}

\subsection{Module 2.1: Exponential Growth and Decay}

\begin{itemize}
    \concept{Discrete-time population growth}
    \concept{Exponential growth}
    \concept{Exponential decay}
    \concept{Recurrence equations}
    \concept{Graphical visualization of recurrence equations}
\end{itemize}

\subsection{Module 2.2: Sequences}

\begin{itemize}
    \concept{Definition of sequences}
    \concept{Recursive sequences}
    \concept{Limits of sequences}
    \concept{Recurrence equations}
    \concept{Summation notation}
\end{itemize}

\subsection{Module 2.3: Modeling with Recurrence Equations}

\begin{itemize}
    \concept{Density-dependent population growth}
    \concept{Beverton--Holt model}
    \concept{Discrete logistic equation}
    \concept{Biological modeling}
\end{itemize}

\subsection{Learning Objectives}

\begin{enumerate}
    \objective{Represent biological growth using discrete-time models.}
    \objective{Construct and interpret recurrence equations.}
    \objective{Generate and interpret sequences.}
    \objective{Analyze limiting behavior of sequences.}
    \objective{Use summation notation.}
    \objective{Interpret density-dependent population models.}
    \objective{Compare exponential and logistic-type growth.}
\end{enumerate}

\subsection{Prerequisites}

\begin{itemize}
    \prereq{Algebra}
    \prereq{Exponential functions}
    \prereq{Function notation}
\end{itemize}

\subsection{Difficult Topics}

\begin{itemize}
    \difficult{Recursive versus explicit representations}
    \difficult{Limits of sequences}
    \difficult{Interpreting recurrence equations dynamically}
    \difficult{Distinguishing exponential and density-dependent growth}
\end{itemize}

% ============================================================
% CHAPTER 3
% ============================================================

\section{Chapter 3: Limits and Continuity}

\subsection{Major Topic}

\textbf{The conceptual foundation of differential calculus.}

\subsection{Module 3.1: Limits}

\begin{itemize}
    \concept{Informal definition of a limit}
    \concept{Limit notation}
    \concept{Limit laws}
    \concept{Pitfalls in calculating limits}
\end{itemize}

\subsection{Module 3.2: Continuity}

\begin{itemize}
    \concept{Continuity}
    \concept{Combinations of continuous functions}
    \concept{Relationship between limits and continuity}
\end{itemize}

\subsection{Module 3.3: Limits at Infinity}

\begin{itemize}
    \concept{Behavior of functions as $x\to\infty$}
    \concept{Behavior of functions as $x\to-\infty$}
    \concept{Horizontal asymptotic behavior}
\end{itemize}

\subsection{Module 3.4: Trigonometric Limits}

\begin{itemize}
    \concept{Fundamental trigonometric limits}
    \concept{Geometric argument for trigonometric limits}
    \concept{Sandwich theorem}
\end{itemize}

\subsection{Module 3.5: Properties of Continuous Functions}

\begin{itemize}
    \concept{Intermediate Value Theorem}
    \concept{Bisection method}
    \concept{Numerical approximation of roots}
\end{itemize}

\subsection{Module 3.6: Formal Definition of Limits}

\textit{Optional/enrichment topic.}

\begin{itemize}
    \concept{Formal limit definition}
    \concept{$\epsilon$--$\delta$ reasoning}
    \concept{Rigorous limit arguments}
\end{itemize}

\subsection{Learning Objectives}

\begin{enumerate}
    \objective{Explain the intuitive meaning of a limit.}
    \objective{Calculate limits using algebraic and limit-law methods.}
    \objective{Identify common errors in limit calculations.}
    \objective{Determine whether a function is continuous.}
    \objective{Analyze limits at infinity.}
    \objective{Use important trigonometric limits.}
    \objective{Apply the Intermediate Value Theorem.}
    \objective{Use bisection to approximate solutions.}
    \objective{Explain the formal definition of a limit when appropriate.}
\end{enumerate}

\subsection{Prerequisites}

\begin{itemize}
    \prereq{Functions and graphing}
    \prereq{Algebraic manipulation}
    \prereq{Trigonometry}
\end{itemize}

\subsection{Difficult Topics}

\begin{itemize}
    \difficult{Understanding a limit as a process rather than a value}
    \difficult{Indeterminate forms}
    \difficult{Limits at infinity}
    \difficult{Trigonometric limits}
    \difficult{Intermediate Value Theorem}
    \difficult{Formal $\epsilon$--$\delta$ definition}
\end{itemize}

% ============================================================
% CHAPTER 4
% ============================================================

\section{Chapter 4: Differentiation}

\subsection{Major Topic}

\textbf{Derivatives, rates of change, and differentiation rules.}

\subsection{Module 4.1: Definition of the Derivative}

\begin{itemize}
    \concept{Derivative from first principles}
    \concept{Difference quotients}
    \concept{Instantaneous rate of change}
    \concept{Tangent lines}
\end{itemize}

\subsection{Module 4.2: Properties of the Derivative}

\begin{itemize}
    \concept{Interpretation of derivatives}
    \concept{Differentiability}
    \concept{Relationship between differentiability and continuity}
\end{itemize}

\subsection{Module 4.3: Basic Differentiation Rules}

\begin{itemize}
    \concept{Power rule}
    \concept{Constant rule}
    \concept{Sum and difference rules}
    \concept{Derivatives of polynomial functions}
\end{itemize}

\subsection{Module 4.4: Product, Quotient, and Power Functions}

\begin{itemize}
    \concept{Product rule}
    \concept{Quotient rule}
    \concept{Derivatives of rational functions}
    \concept{Derivatives of power functions}
\end{itemize}

\subsection{Module 4.5: Chain Rule}

\begin{itemize}
    \concept{Composition of functions}
    \concept{Chain rule}
    \concept{Nested functions}
\end{itemize}

\subsection{Module 4.6: Implicit Differentiation and Related Rates}

\begin{itemize}
    \concept{Implicit functions}
    \concept{Implicit differentiation}
    \concept{Related rates}
\end{itemize}

\subsection{Module 4.7: Higher Derivatives}

\begin{itemize}
    \concept{Second derivative}
    \concept{Higher-order derivatives}
\end{itemize}

\subsection{Module 4.8: Trigonometric Derivatives}

\begin{itemize}
    \concept{Derivatives of trigonometric functions}
\end{itemize}

\subsection{Module 4.9: Exponential Derivatives}

\begin{itemize}
    \concept{Derivative of exponential functions}
\end{itemize}

\subsection{Module 4.10: Inverse and Logarithmic Derivatives}

\begin{itemize}
    \concept{Derivative of inverse functions}
    \concept{Derivative of logarithmic functions}
    \concept{Derivative of inverse tangent}
    \concept{Logarithmic differentiation}
\end{itemize}

\subsection{Module 4.11: Linear Approximation and Error Propagation}

\begin{itemize}
    \concept{Linearization}
    \concept{Approximation}
    \concept{Differential estimates}
    \concept{Error propagation}
\end{itemize}

\subsection{Learning Objectives}

\begin{enumerate}
    \objective{Define the derivative using a limit.}
    \objective{Interpret the derivative as an instantaneous rate of change.}
    \objective{Interpret the derivative geometrically as a tangent slope.}
    \objective{Differentiate elementary functions.}
    \objective{Apply the product, quotient, and chain rules.}
    \objective{Differentiate implicitly defined functions.}
    \objective{Solve related-rate problems.}
    \objective{Compute and interpret higher derivatives.}
    \objective{Differentiate trigonometric, exponential, logarithmic,
    and inverse functions.}
    \objective{Use linear approximation to estimate function values.}
    \objective{Estimate and interpret propagated errors.}
\end{enumerate}

\subsection{Prerequisites}

\begin{itemize}
    \prereq{Chapter 3: Limits and Continuity}
    \prereq{Elementary functions}
    \prereq{Algebra and trigonometry}
\end{itemize}

\subsection{Difficult Topics}

\begin{itemize}
    \difficult{Derivative from first principles}
    \difficult{Chain rule}
    \difficult{Implicit differentiation}
    \difficult{Related rates}
    \difficult{Inverse-function derivatives}
    \difficult{Logarithmic differentiation}
    \difficult{Error propagation}
\end{itemize}

% ============================================================
% CHAPTER 5
% ============================================================

\section{Chapter 5: Applications of Differentiation}

\subsection{Major Topic}

\textbf{Using derivatives to understand, analyze, and optimize models.}

\subsection{Module 5.1: Extrema and the Mean Value Theorem}

\begin{itemize}
    \concept{Extreme Value Theorem}
    \concept{Local extrema}
    \concept{Mean Value Theorem}
\end{itemize}

\subsection{Module 5.2: Monotonicity and Concavity}

\begin{itemize}
    \concept{Increasing and decreasing functions}
    \concept{First derivative}
    \concept{Concavity}
    \concept{Second derivative}
\end{itemize}

\subsection{Module 5.3: Extrema and Inflection Points}

\begin{itemize}
    \concept{Critical points}
    \concept{Local and global extrema}
    \concept{Inflection points}
\end{itemize}

\subsection{Module 5.4: Optimization}

\begin{itemize}
    \concept{Optimization models}
    \concept{Objective functions}
    \concept{Constraints}
    \concept{Critical-point analysis}
\end{itemize}

\subsection{Module 5.5: L'Hopital's Rule}

\begin{itemize}
    \concept{Indeterminate forms}
    \concept{Limit evaluation using derivatives}
\end{itemize}

\subsection{Module 5.6: Graphing and Asymptotes}

\begin{itemize}
    \concept{Derivative-based graphing}
    \concept{Asymptotic behavior}
    \concept{Function shape}
\end{itemize}

\subsection{Optional Modules}

\begin{itemize}
    \concept{Stability of recurrence equations}
    \concept{Newton--Raphson numerical method}
    \concept{Biological systems modeled using differential equations}
\end{itemize}

\subsection{Module 5.10: Antiderivatives}

\begin{itemize}
    \concept{Antiderivatives}
    \concept{Reverse differentiation}
    \concept{Indefinite integration as preparation for Chapter 6}
\end{itemize}

\subsection{Learning Objectives}

\begin{enumerate}
    \objective{Identify and classify critical points.}
    \objective{Determine intervals of increase and decrease.}
    \objective{Determine concavity and inflection points.}
    \objective{Use the Extreme Value and Mean Value Theorems.}
    \objective{Solve applied optimization problems.}
    \objective{Evaluate appropriate indeterminate limits using
    L'Hopital's Rule.}
    \objective{Use derivatives to sketch and analyze graphs.}
    \objective{Find antiderivatives of basic functions.}
\end{enumerate}

\subsection{Prerequisites}

\begin{itemize}
    \prereq{Chapter 4: Differentiation}
    \prereq{Limits and continuity}
    \prereq{Algebraic modeling}
\end{itemize}

\subsection{Difficult Topics}

\begin{itemize}
    \difficult{Distinguishing local and global extrema}
    \difficult{Optimization modeling}
    \difficult{Concavity and inflection points}
    \difficult{Graphing from derivative information}
    \difficult{Choosing an appropriate derivative-based method}
\end{itemize}

% ============================================================
% CHAPTER 6
% ============================================================

\section{Chapter 6: Integration}

\subsection{Major Topic}

\textbf{Accumulation, area, antiderivatives, and the Fundamental
Theorem of Calculus.}

\subsection{Module 6.1: Definite Integral}

\begin{itemize}
    \concept{Area problem}
    \concept{Riemann integral}
    \concept{Riemann sums}
    \concept{Properties of the definite integral}
\end{itemize}

\subsection{Module 6.2: Fundamental Theorem of Calculus}

\begin{itemize}
    \concept{FTC Part I}
    \concept{Antiderivatives}
    \concept{Indefinite integrals}
    \concept{FTC Part II}
    \concept{Relationship between differentiation and integration}
\end{itemize}

\subsection{Module 6.3: Applications of Integration}

\begin{itemize}
    \concept{Cumulative change}
    \concept{Average values}
    \concept{Areas}
    \concept{Volumes}
    \concept{Arc length}
\end{itemize}

\subsection{Learning Objectives}

\begin{enumerate}
    \objective{Explain the definite integral as an accumulation process.}
    \objective{Approximate area using Riemann sums.}
    \objective{Evaluate definite integrals.}
    \objective{Distinguish definite from indefinite integrals.}
    \objective{Use both parts of the Fundamental Theorem of Calculus.}
    \objective{Interpret integrals as cumulative change.}
    \objective{Calculate and interpret average values.}
    \objective{Apply integration to areas and volumes when appropriate.}
\end{enumerate}

\subsection{Prerequisites}

\begin{itemize}
    \prereq{Limits}
    \prereq{Continuity}
    \prereq{Differentiation}
    \prereq{Antiderivatives}
\end{itemize}

\subsection{Difficult Topics}

\begin{itemize}
    \difficult{Riemann sums}
    \difficult{Understanding accumulation}
    \difficult{The Fundamental Theorem of Calculus}
    \difficult{Distinguishing rate from accumulated quantity}
    \difficult{Setting up applied integrals}
\end{itemize}

% ============================================================
% CHAPTER 7
% ============================================================

\section{Chapter 7: Integration Techniques and Computational Methods}

\subsection{Major Topic}

\textbf{Symbolic and numerical methods for integration.}

\subsection{Module 7.1: Substitution Rule}

\begin{itemize}
    \concept{$u$-substitution}
    \concept{Indefinite integrals}
    \concept{Definite integrals}
\end{itemize}

\subsection{Module 7.2: Integration by Parts}

\begin{itemize}
    \concept{Integration by parts}
    \concept{Selection of factors}
    \concept{Repeated integration}
\end{itemize}

\subsection{Module 7.3: Rational Functions and Partial Fractions}

\begin{itemize}
    \concept{Proper rational functions}
    \concept{Partial-fraction decomposition}
    \concept{Repeated linear factors}
    \concept{Irreducible quadratic factors}
\end{itemize}

\subsection{Module 7.4: Improper Integrals}

\textit{Optional/enrichment topic.}

\begin{itemize}
    \concept{Unbounded intervals}
    \concept{Unbounded integrands}
    \concept{Comparison methods}
\end{itemize}

\subsection{Module 7.5: Numerical Integration}

\begin{itemize}
    \concept{Midpoint rule}
    \concept{Trapezoidal rule}
    \concept{Spreadsheet-based numerical integration}
    \concept{Numerical error estimation}
\end{itemize}

\subsection{Module 7.6: Taylor Approximation}

\textit{Optional/enrichment topic.}

\begin{itemize}
    \concept{Taylor polynomials}
    \concept{Taylor polynomial about $x=a$}
    \concept{Approximation accuracy}
\end{itemize}

\subsection{Learning Objectives}

\begin{enumerate}
    \objective{Use substitution to simplify integrals.}
    \objective{Apply integration by parts.}
    \objective{Integrate selected rational functions using partial fractions.}
    \objective{Recognize when numerical integration is appropriate.}
    \objective{Apply midpoint and trapezoidal rules.}
    \objective{Interpret numerical approximation error.}
    \objective{Use Taylor polynomials as local approximations
    when appropriate.}
\end{enumerate}

\subsection{Prerequisites}

\begin{itemize}
    \prereq{Chapter 6: Integration}
    \prereq{Differentiation rules}
    \prereq{Algebraic manipulation}
\end{itemize}

\subsection{Difficult Topics}

\begin{itemize}
    \difficult{Recognizing the correct substitution}
    \difficult{Integration by parts}
    \difficult{Partial-fraction decomposition}
    \difficult{Improper integrals}
    \difficult{Error analysis in numerical integration}
\end{itemize}

% ============================================================
% CHAPTER 8
% ============================================================

\section{Chapter 8: Differential Equations}

\subsection{Major Topic}

\textbf{Using calculus to model dynamic biological and social systems.}

\subsection{Module 8.1: Separable Differential Equations}

\begin{itemize}
    \concept{Pure-time differential equations}
    \concept{Autonomous differential equations}
    \concept{Separable equations}
    \concept{Initial-value interpretation}
\end{itemize}

\subsection{Module 8.2: Equilibria and Stability}

\begin{itemize}
    \concept{Equilibrium points}
    \concept{Graphical equilibrium analysis}
    \concept{Stability}
    \concept{Vector fields}
    \concept{Behavior near equilibria}
\end{itemize}

\subsection{Module 8.3: Differential Equation Models}

\begin{itemize}
    \concept{Compartment models}
    \concept{Ecological models}
    \concept{Chemical reaction models}
    \concept{Models of cooperation}
    \concept{Epidemic models}
\end{itemize}

\subsection{Module 8.4: Integrating Factors and Two-Compartment Models}

\begin{itemize}
    \concept{Integrating factors}
    \concept{Two-compartment systems}
\end{itemize}

\subsection{Learning Objectives}

\begin{enumerate}
    \objective{Explain what a differential equation represents.}
    \objective{Solve basic separable differential equations.}
    \objective{Identify equilibrium solutions.}
    \objective{Determine qualitative stability of equilibria.}
    \objective{Interpret vector-field representations.}
    \objective{Construct compartment and ecological models.}
    \objective{Apply differential equations to chemical and epidemic models.}
    \objective{Use integrating factors for appropriate first-order equations.}
\end{enumerate}

\subsection{Prerequisites}

\begin{itemize}
    \prereq{Differentiation}
    \prereq{Integration}
    \prereq{Algebra}
    \prereq{Function interpretation}
\end{itemize}

\subsection{Difficult Topics}

\begin{itemize}
    \difficult{Translating a verbal model into a differential equation}
    \difficult{Separation of variables}
    \difficult{Equilibrium and stability analysis}
    \difficult{Interpreting vector fields}
    \difficult{Compartment modeling}
\end{itemize}

% ============================================================
% CHAPTER 9
% ============================================================

\section{Chapter 9: Linear Algebra and Analytic Geometry}

\subsection{Major Topic}

\textbf{Linear systems, matrices, eigenvalues, and vectors.}

\subsection{Module 9.1: Linear Systems}

\begin{itemize}
    \concept{Graphical solution}
    \concept{Elimination}
    \concept{Systems of equations}
    \concept{Matrix representation}
\end{itemize}

\subsection{Module 9.2: Matrices}

\begin{itemize}
    \concept{Matrix operations}
    \concept{Matrix multiplication}
    \concept{Inverse matrices}
\end{itemize}

\subsection{Module 9.3: Linear Maps, Eigenvectors, and Eigenvalues}

\begin{itemize}
    \concept{Linear transformations}
    \concept{Eigenvectors}
    \concept{Eigenvalues}
    \concept{Graphical representation}
\end{itemize}

\subsection{Module 9.4: Demographic Modeling}

\textit{Optional/enrichment topic.}

\begin{itemize}
    \concept{Leslie matrices}
    \concept{Stable age distributions}
    \concept{Population dynamics}
\end{itemize}

\subsection{Module 9.5: Analytic Geometry}

\begin{itemize}
    \concept{Points and vectors in higher dimensions}
    \concept{Dot product}
    \concept{Parametric equations of lines}
\end{itemize}

\subsection{Learning Objectives}

\begin{enumerate}
    \objective{Solve systems of linear equations.}
    \objective{Represent linear systems using matrices.}
    \objective{Perform matrix operations.}
    \objective{Interpret matrix multiplication.}
    \objective{Understand inverse matrices.}
    \objective{Identify eigenvalues and eigenvectors.}
    \objective{Interpret vectors geometrically.}
    \objective{Use dot products and parametric equations.}
\end{enumerate}

\subsection{Prerequisites}

\begin{itemize}
    \prereq{Algebra}
    \prereq{Systems of equations}
    \prereq{Basic coordinate geometry}
\end{itemize}

\subsection{Difficult Topics}

\begin{itemize}
    \difficult{Matrix multiplication}
    \difficult{Inverse matrices}
    \difficult{Eigenvalues and eigenvectors}
    \difficult{Transition from scalar to vector thinking}
\end{itemize}

% ============================================================
% CHAPTER 10
% ============================================================

\section{Chapter 10: Multivariable Calculus}

\subsection{Major Topic}

\textbf{Functions, derivatives, and optimization involving multiple
independent variables.}

\subsection{Module 10.1: Functions of Several Variables}

\begin{itemize}
    \concept{Functions of two variables}
    \concept{Functions of more than two variables}
    \concept{Surface plots}
    \concept{Heat maps}
    \concept{Contour plots}
\end{itemize}

\subsection{Module 10.2: Limits and Continuity}

\textit{Optional/enrichment topic.}

\begin{itemize}
    \concept{Multivariable limits}
    \concept{Continuity}
    \concept{Formal limits}
\end{itemize}

\subsection{Module 10.3: Partial Derivatives}

\begin{itemize}
    \concept{Partial derivatives}
    \concept{Higher-order partial derivatives}
\end{itemize}

\subsection{Module 10.4: Tangent Planes, Differentiability, and Linearization}

\begin{itemize}
    \concept{Tangent planes}
    \concept{Differentiability}
    \concept{Multivariable linearization}
    \concept{Vector-valued functions}
\end{itemize}

\subsection{Module 10.5: Multivariable Chain Rule}

\textit{Optional/enrichment topic.}

\begin{itemize}
    \concept{Chain rule for functions of several variables}
    \concept{Implicit differentiation}
\end{itemize}

\subsection{Module 10.6: Directional Derivatives and Gradient}

\textit{Optional/enrichment topic.}

\begin{itemize}
    \concept{Directional derivative}
    \concept{Gradient vector}
    \concept{Properties of the gradient}
\end{itemize}

\subsection{Module 10.7: Multivariable Optimization}

\textit{Optional/enrichment topic.}

\begin{itemize}
    \concept{Local maxima and minima}
    \concept{Global extrema}
    \concept{Constrained extrema}
    \concept{Least-squares data fitting}
\end{itemize}

\subsection{Module 10.8--10.9: Mathematical Models}

\begin{itemize}
    \concept{Diffusion}
    \concept{Systems of recurrence equations}
    \concept{Equilibria and stability}
\end{itemize}

\subsection{Learning Objectives}

\begin{enumerate}
    \objective{Represent functions of multiple variables.}
    \objective{Interpret surface, heat-map, and contour representations.}
    \objective{Calculate partial derivatives.}
    \objective{Interpret higher-order partial derivatives.}
    \objective{Construct tangent-plane and linear approximations.}
    \objective{Apply the multivariable chain rule when appropriate.}
    \objective{Interpret directional derivatives and gradients.}
    \objective{Solve basic multivariable optimization problems.}
    \objective{Apply multivariable concepts to mathematical models.}
\end{enumerate}

\subsection{Prerequisites}

\begin{itemize}
    \prereq{Single-variable differentiation}
    \prereq{Single-variable optimization}
    \prereq{Algebra}
    \prereq{Vectors and coordinate geometry}
\end{itemize}

\subsection{Difficult Topics}

\begin{itemize}
    \difficult{Visualizing functions of two variables}
    \difficult{Partial derivatives}
    \difficult{Tangent planes}
    \difficult{Multivariable chain rule}
    \difficult{Gradient vectors}
    \difficult{Constrained optimization}
\end{itemize}

% ============================================================
% CHAPTER 11
% ============================================================

\section{Chapter 11: Systems of Differential Equations}

\subsection{Major Topic}

\textbf{Dynamic systems involving multiple interacting variables.}

\subsection{Module 11.1: Linear Systems}

\begin{itemize}
    \concept{Vector fields}
    \concept{Solutions of linear systems}
    \concept{Equilibria}
    \concept{Stability}
    \concept{Complex conjugate eigenvalues}
\end{itemize}

\subsection{Module 11.2: Applications}

\begin{itemize}
    \concept{Two-compartment models}
    \concept{Interaction models}
    \concept{Harmonic oscillator}
\end{itemize}

\subsection{Module 11.3: Nonlinear Autonomous Systems}

\begin{itemize}
    \concept{Analytical approaches}
    \concept{Graphical analysis}
    \concept{$2\times2$ nonlinear systems}
\end{itemize}

\subsection{Module 11.4: Lotka--Volterra Models}

\begin{itemize}
    \concept{Competition}
    \concept{Predator--prey interactions}
    \concept{Interspecific interactions}
\end{itemize}

\subsection{Module 11.5: Additional Mathematical Models}

\textit{Optional/enrichment topic.}

\begin{itemize}
    \concept{Community matrix}
    \concept{Neuron activity}
    \concept{Enzymatic reactions}
    \concept{Microbial growth}
    \concept{Epidemic models}
\end{itemize}

\subsection{Learning Objectives}

\begin{enumerate}
    \objective{Represent interacting dynamic variables using systems
    of differential equations.}
    \objective{Interpret vector fields.}
    \objective{Analyze equilibrium points and stability.}
    \objective{Understand the role of eigenvalues in linear systems.}
    \objective{Analyze nonlinear two-dimensional systems graphically.}
    \objective{Apply systems to ecological and biological models.}
    \objective{Interpret predator--prey and competition dynamics.}
\end{enumerate}

\subsection{Prerequisites}

\begin{itemize}
    \prereq{Chapter 8: Differential Equations}
    \prereq{Chapter 9: Linear Algebra}
    \prereq{Eigenvalues and eigenvectors}
\end{itemize}

\subsection{Difficult Topics}

\begin{itemize}
    \difficult{Connecting matrices to differential equations}
    \difficult{Eigenvalue-based stability analysis}
    \difficult{Phase-plane reasoning}
    \difficult{Complex eigenvalues}
    \difficult{Nonlinear systems}
\end{itemize}

% ============================================================
% CHAPTER 12
% ============================================================

\section{Chapter 12: Probability and Statistics}

\subsection{Major Topic}

\textbf{Probability, random variables, distributions, limit theorems,
and statistical tools.}

\subsection{Module 12.1: Counting}

\begin{itemize}
    \concept{Multiplication principle}
    \concept{Permutations}
    \concept{Combinations}
    \concept{Combined counting principles}
\end{itemize}

\subsection{Module 12.2: Probability}

\begin{itemize}
    \concept{Basic probability definitions}
    \concept{Equally likely outcomes}
\end{itemize}

\subsection{Module 12.3: Conditional Probability and Independence}

\begin{itemize}
    \concept{Conditional probability}
    \concept{Law of total probability}
    \concept{Independence}
    \concept{Bayes' formula}
\end{itemize}

\subsection{Module 12.4: Discrete Random Variables}

\begin{itemize}
    \concept{Discrete distributions}
    \concept{Mean}
    \concept{Variance}
    \concept{Binomial distribution}
    \concept{Multinomial distribution}
    \concept{Geometric distribution}
    \concept{Poisson distribution}
\end{itemize}

\subsection{Module 12.5: Continuous Distributions}

\begin{itemize}
    \concept{Density functions}
    \concept{Normal distribution}
    \concept{Uniform distribution}
    \concept{Exponential distribution}
    \concept{Poisson process}
    \concept{Aging}
\end{itemize}

\subsection{Module 12.6: Limit Theorems}

\begin{itemize}
    \concept{Law of Large Numbers}
    \concept{Central Limit Theorem}
\end{itemize}

\subsection{Module 12.7: Statistical Tools}

\begin{itemize}
    \concept{Describing univariate data}
    \concept{Parameter estimation}
    \concept{Linear regression}
\end{itemize}

\subsection{Learning Objectives}

\begin{enumerate}
    \objective{Apply basic counting principles.}
    \objective{Calculate elementary probabilities.}
    \objective{Use conditional probability and independence.}
    \objective{Apply Bayes' formula.}
    \objective{Calculate means and variances of discrete random variables.}
    \objective{Recognize common discrete and continuous distributions.}
    \objective{Interpret probability density functions.}
    \objective{Explain the Law of Large Numbers and Central Limit Theorem.}
    \objective{Summarize univariate data.}
    \objective{Estimate parameters from data.}
    \objective{Interpret linear regression models.}
\end{enumerate}

\subsection{Prerequisites}

\begin{itemize}
    \prereq{Algebra}
    \prereq{Function notation}
    \prereq{Basic summation}
    \prereq{For continuous distributions: basic integration}
\end{itemize}

\subsection{Difficult Topics}

\begin{itemize}
    \difficult{Conditional probability}
    \difficult{Independence}
    \difficult{Bayes' formula}
    \difficult{Expected value and variance}
    \difficult{Continuous probability density functions}
    \difficult{Central Limit Theorem}
    \difficult{Parameter estimation}
\end{itemize}

% ============================================================
% 3. CROSS-COURSE CONCEPT MAP
% ============================================================

\section{Cross-Course Concept Map}

The major mathematical ideas can be organized into the following
conceptual hierarchy.

\subsection{Representation}

\[
\text{Verbal Model}
\leftrightarrow
\text{Equation}
\leftrightarrow
\text{Table}
\leftrightarrow
\text{Graph}
\]

Students repeatedly move among these representations.

\subsection{Change}

\[
\text{Average Rate of Change}
\rightarrow
\text{Instantaneous Rate of Change}
\rightarrow
\text{Derivative}
\]

\subsection{Accumulation}

\[
\text{Small Contributions}
\rightarrow
\text{Riemann Sum}
\rightarrow
\text{Definite Integral}
\rightarrow
\text{Cumulative Quantity}
\]

\subsection{Derivative--Integral Relationship}

\[
\boxed{
\text{Differentiation}
\quad\Longleftrightarrow\quad
\text{Integration}
}
\]

through the Fundamental Theorem of Calculus.

\subsection{Dynamic Modeling}

\[
\text{State}
\rightarrow
\text{Rate of Change}
\rightarrow
\text{Differential Equation}
\rightarrow
\text{Solution}
\rightarrow
\text{Prediction}
\]

\subsection{Systems Thinking}

\[
\text{Several Variables}
\rightarrow
\text{Linear Algebra}
\rightarrow
\text{Systems}
\rightarrow
\text{Eigenvalues}
\rightarrow
\text{Stability}.
\]

% ============================================================
% 4. DIFFICULTY PROGRESSION
% ============================================================

\section{Difficulty Progression}

\begin{longtable}{p{0.17\textwidth} p{0.28\textwidth} p{0.45\textwidth}}
\toprule
\textbf{Level} & \textbf{Topics} & \textbf{Main Challenge}\\
\midrule
Foundational &
Algebra, functions, graphing &
Fluency with symbolic and graphical representations\\
\addlinespace
Developing &
Sequences, limits, continuity &
Understanding mathematical processes and limiting behavior\\
\addlinespace
Intermediate &
Differentiation &
Connecting formulas with rates of change and geometric meaning\\
\addlinespace
Intermediate--Advanced &
Applications of derivatives &
Translating real-world problems into mathematical optimization models\\
\addlinespace
Advanced &
Integration &
Understanding accumulation and the Fundamental Theorem\\
\addlinespace
Advanced &
Integration techniques &
Selecting an appropriate symbolic or numerical method\\
\addlinespace
Advanced &
Differential equations &
Connecting rates of change to dynamic models\\
\addlinespace
Advanced &
Linear algebra &
Moving from scalar equations to vector/matrix representations\\
\addlinespace
Advanced &
Multivariable calculus &
Generalizing one-variable ideas to several dimensions\\
\addlinespace
Advanced &
Systems of differential equations &
Combining calculus, matrices, eigenvalues, and dynamical systems\\
\addlinespace
Application-focused &
Probability and statistics &
Connecting mathematical distributions to uncertainty and data\\
\bottomrule
\end{longtable}

% ============================================================
% 5. HIGH-PRIORITY DIFFICULT TOPICS
% ============================================================

\section{High-Priority Difficult Topics}

The following topics deserve additional instructional time because
they require substantial conceptual transitions.

\begin{enumerate}
    \item \textbf{Function interpretation}
    
    Students must move fluently among equations, graphs, tables,
    and verbal descriptions.
    
    \item \textbf{Limits}
    
    The idea of approaching a value is conceptually different from
    ordinary algebraic substitution.
    
    \item \textbf{Derivative as instantaneous rate of change}
    
    Students must connect the geometric, algebraic, and applied
    meanings of the derivative.
    
    \item \textbf{Chain Rule}
    
    Correctly identifying nested functions and their dependencies
    is a major symbolic challenge.
    
    \item \textbf{Related Rates}
    
    Students must translate a changing real-world situation into
    equations involving derivatives.
    
    \item \textbf{Optimization}
    
    The main challenge is usually model formulation rather than
    differentiation itself.
    
    \item \textbf{Fundamental Theorem of Calculus}
    
    Students must understand why differentiation and integration
    are connected.
    
    \item \textbf{Setting up integrals}
    
    Applied integration requires identifying the quantity being
    accumulated before performing calculations.
    
    \item \textbf{Differential-equation modeling}
    
    Students must translate biological or social mechanisms into
    mathematical rate equations.
    
    \item \textbf{Equilibrium and stability}
    
    Students must reason dynamically rather than only calculate
    numerical answers.
    
    \item \textbf{Eigenvalues and eigenvectors}
    
    These require a transition from elementary algebra to structural
    linear-algebra reasoning.
    
    \item \textbf{Multivariable visualization}
    
    Students must reason about surfaces, contours, and changes in
    multiple directions.
\end{enumerate}

% ============================================================
% 6. RECOMMENDED COURSE MODULES
% ============================================================

\section{Recommended Module-Based Course Structure}

For instructional purposes, the twelve textbook chapters can be
grouped into the following larger modules.

\begin{longtable}{p{0.12\textwidth} p{0.28\textwidth} p{0.48\textwidth}}
\toprule
\textbf{Module} & \textbf{Theme} & \textbf{Chapters/Topics}\\
\midrule
M1 &
Mathematical Foundations &
Real numbers, algebra, trigonometry, functions, graphing\\
\addlinespace
M2 &
Discrete Modeling &
Exponential growth, sequences, recurrence equations, population models\\
\addlinespace
M3 &
Limits and Continuity &
Limits, continuity, limits at infinity, IVT\\
\addlinespace
M4 &
Differential Calculus &
Derivative definition, rules, chain rule, implicit differentiation,
related rates\\
\addlinespace
M5 &
Applications of Derivatives &
Extrema, monotonicity, concavity, optimization, graphing\\
\addlinespace
M6 &
Integral Calculus &
Definite integrals, Riemann sums, FTC, accumulation\\
\addlinespace
M7 &
Integration Methods &
Substitution, integration by parts, partial fractions,
numerical integration\\
\addlinespace
M8 &
Differential Equations &
Separable equations, equilibria, stability, biological models\\
\addlinespace
M9 &
Linear Algebra &
Linear systems, matrices, eigenvalues, eigenvectors\\
\addlinespace
M10 &
Multivariable Calculus &
Multivariable functions, partial derivatives, tangent planes,
gradients, optimization\\
\addlinespace
M11 &
Dynamical Systems &
Systems of differential equations, stability, ecological models\\
\addlinespace
M12 &
Probability and Statistics &
Probability, random variables, distributions, limit theorems,
regression\\
\bottomrule
\end{longtable}

% ============================================================
% 7. CORE VS OPTIONAL CONTENT
% ============================================================

\section{Core Versus Optional Content}

The textbook explicitly indicates that asterisked sections are not
required for progressing through the later sections of the text.
Therefore, a practical course can distinguish between core and
enrichment material.

\subsection{Core Calculus Sequence}

\begin{enumerate}
    \item Chapter 1: Preview and Review
    \item Chapter 2: Discrete-Time Models, Sequences, and Difference
    Equations
    \item Chapter 3: Limits and Continuity
    \item Chapter 4: Differentiation
    \item Chapter 5: Applications of Differentiation
    \item Chapter 6: Integration
    \item Chapter 7: Integration Techniques and Computational Methods
    \item Chapter 8: Differential Equations
\end{enumerate}

\subsection{Advanced/Extension Sequence}

\begin{enumerate}
    \item Chapter 9: Linear Algebra and Analytic Geometry
    \item Chapter 10: Multivariable Calculus
    \item Chapter 11: Systems of Differential Equations
    \item Chapter 12: Probability and Statistics
\end{enumerate}

\subsection{Typical Optional/Enrichment Areas}

Examples of topics explicitly marked as optional in the supplied
contents include:

\begin{itemize}
    \item Formal definitions of limits
    \item Spreadsheet implementations
    \item Drug-absorption modeling
    \item Stability of recurrence equations
    \item Newton--Raphson method
    \item Differential-equation modeling extensions
    \item Rigorous proof of the Fundamental Theorem
    \item Improper integrals
    \item Taylor approximation
    \item Multivariable limits
    \item Multivariable chain rule
    \item Directional derivatives and gradients
    \item Multivariable optimization
    \item Diffusion and systems of recurrence equations
    \item Additional differential-equation models
\end{itemize}

% ============================================================
% 8. ASSESSMENT ALIGNMENT
% ============================================================

\section{Suggested Assessment Alignment}

\begin{longtable}{p{0.22\textwidth} p{0.30\textwidth} p{0.35\textwidth}}
\toprule
\textbf{Assessment Type} & \textbf{Primary Skills} & \textbf{Best-Aligned Topics}\\
\midrule
Diagnostic &
Prerequisite fluency &
Algebra, functions, trigonometry\\
\addlinespace
Concept quizzes &
Conceptual understanding &
Limits, continuity, derivatives, integrals\\
\addlinespace
Problem sets &
Symbolic proficiency &
Differentiation, integration, differential equations\\
\addlinespace
Modeling assignments &
Mathematical translation &
Population, epidemic, ecological, compartment models\\
\addlinespace
Graphing tasks &
Visual reasoning &
Functions, derivatives, concavity, vector fields\\
\addlinespace
Application projects &
Interpretation and communication &
Optimization, biological models, dynamical systems\\
\addlinespace
Cumulative examinations &
Integration of concepts &
Limits through integration and differential equations\\
\bottomrule
\end{longtable}

% ============================================================
% 9. LEARNING OUTCOME HIERARCHY
% ============================================================

\section{Learning Outcome Hierarchy}

\subsection{Level 1: Mathematical Fluency}

Students can:

\begin{itemize}
    \item manipulate algebraic expressions;
    \item evaluate functions;
    \item interpret graphs;
    \item use exponential, logarithmic, and trigonometric functions.
\end{itemize}

\subsection{Level 2: Conceptual Calculus}

Students can:

\begin{itemize}
    \item explain limits;
    \item determine continuity;
    \item explain derivatives as rates of change;
    \item explain integrals as accumulation.
\end{itemize}

\subsection{Level 3: Computational Calculus}

Students can:

\begin{itemize}
    \item calculate limits;
    \item differentiate functions;
    \item integrate functions;
    \item use standard calculus rules and techniques.
\end{itemize}

\subsection{Level 4: Analytical Reasoning}

Students can:

\begin{itemize}
    \item analyze extrema;
    \item determine monotonicity and concavity;
    \item optimize functions;
    \item analyze equilibrium and stability.
\end{itemize}

\subsection{Level 5: Mathematical Modeling}

Students can:

\begin{itemize}
    \item formulate mathematical models;
    \item interpret model parameters;
    \item analyze model behavior;
    \item connect mathematical results to biological or social phenomena.
\end{itemize}

\subsection{Level 6: Advanced Mathematical Modeling}

Students can:

\begin{itemize}
    \item work with matrices and eigenvalues;
    \item analyze multivariable functions;
    \item study systems of differential equations;
    \item use probability and statistical models.
\end{itemize}

% ============================================================
% 10. FINAL COURSE MAP
% ============================================================

\section{Final Course Map}

The entire textbook can be understood through five major mathematical
strands.

\begin{enumerate}
    \item \textbf{Functions and Representation}
    
    \[
    \text{Functions}
    \rightarrow
    \text{Graphs}
    \rightarrow
    \text{Models}
    \]
    
    \item \textbf{Change}
    
    \[
    \text{Limits}
    \rightarrow
    \text{Derivatives}
    \rightarrow
    \text{Optimization}
    \]
    
    \item \textbf{Accumulation}
    
    \[
    \text{Riemann Sums}
    \rightarrow
    \text{Integrals}
    \rightarrow
    \text{Cumulative Change}
    \]
    
    \item \textbf{Dynamics}
    
    \[
    \text{Differential Equations}
    \rightarrow
    \text{Equilibria}
    \rightarrow
    \text{Stability}
    \rightarrow
    \text{Systems}
    \]
    
    \item \textbf{Data and Uncertainty}
    
    \[
    \text{Probability}
    \rightarrow
    \text{Distributions}
    \rightarrow
    \text{Limit Theorems}
    \rightarrow
    \text{Statistics}
    \]
\end{enumerate}

\subsection{Complete Dependency Chain}

\[
\boxed{
\begin{array}{c}
\text{Precalculus}
\\
\downarrow
\\
\text{Functions and Graphs}
\\
\downarrow
\\
\text{Discrete Models and Sequences}
\\
\downarrow
\\
\text{Limits and Continuity}
\\
\downarrow
\\
\text{Differentiation}
\\
\downarrow
\\
\text{Applications of Differentiation}
\\
\downarrow
\\
\text{Integration}
\\
\downarrow
\\
\text{Integration Techniques}
\\
\downarrow
\\
\text{Differential Equations}
\\
\downarrow
\\
\text{Linear Algebra}
\\
\downarrow
\\
\text{Multivariable Calculus}
\\
\downarrow
\\
\text{Systems of Differential Equations}
\end{array}
}
\]

Probability and statistics form a complementary application strand:

\[
\boxed{
\text{Counting}
\rightarrow
\text{Probability}
\rightarrow
\text{Random Variables}
\rightarrow
\text{Distributions}
\rightarrow
\text{Limit Theorems}
\rightarrow
\text{Statistics}.
}
\]

% ============================================================
% CONCLUSION
% ============================================================

\section{Conclusion}

The textbook supports a course architecture that begins with
precalculus and function fluency, develops the central calculus
sequence of limits, derivatives, and integrals, and then applies
these tools to biological and social-science models.

The most important conceptual spine is

\[
\boxed{
\text{Functions}
\rightarrow
\text{Limits}
\rightarrow
\text{Derivatives}
\rightarrow
\text{Optimization}
\rightarrow
\text{Integrals}
\rightarrow
\text{Differential Equations}.
}
\]

The later chapters broaden the mathematical toolkit through linear
algebra, multivariable calculus, systems of differential equations,
and probability/statistics.

For a conventional ``Calculus I'' course, the highest-priority
material would normally be Chapters 1--6, with selected material
from Chapters 7--8 used for computational and biological modeling
applications. Chapters 9--12 are best viewed as extensions of the
calculus foundation rather than as the minimum core of a traditional
one-semester Calculus I course.

\end{document}